%-------------------------------------------------------------------------------
%	SECTION TITLE
%-------------------------------------------------------------------------------
\cvsection{Proyectos Relevantes}


%-------------------------------------------------------------------------------
%	CONTENT
%-------------------------------------------------------------------------------
\begin{cventries}

%---------------------------------------------------------
  \cventry
    {}
    {Desarrollo de un pipeline MLOps para la predicción de enfermedades cardiovasculares}
    {}
    {Jul 2025 - Aug 2025}
    {
      \begin{cvitems}
        \item Desarrollo de un pipeline MLOps para predecir enfermedades cardiovasculares a partir de datos clínicos y demográficos.
        \item Modelo KNN alcanzó una exactitud de 0.89 y F1 de 0.89, mejorando frente a modelos base.
        \item Despliegue como API REST en Google Cloud Run con monitoreo de drift y alertas en Slack.
      \end{cvitems}
      \vspace{14pt}
      \begin{itemize}[leftmargin=*]
        \small
        \item[] \textbf{Tecnologías}: MLflow, Prefect, FastAPI, Docker, Terraform, GCP, Evidently AI.
        \hspace{1em}
        \textbf{Enlaces}: 
        \href{https://github.com/aletbm/Cardiovascular-Diseases_E2E_MLPipeline}{Repositorio} | 
        \href{https://github.com/aletbm/Cardiovascular-Diseases_EDA-Modeling/blob/184a432357f98b17ef77ee52ce4994e544a8fc82/cardiovascular-diseases-eda-modeling.ipynb}{Análisis} | 
        \href{https://cardiovascular-diseases-api-761922006747.us-east1.run.app}{API} | 
        \href{https://aletbm.github.io/Cardiovascular-Diseases_E2E_MLPipeline/monitoring/full_monitor_report.html}{Reporte}
      \end{itemize}
    }

%---------------------------------------------------------
  \cventry
    {}
    {Clasificación y segmentación de células sanguíneas cancerígenas – Proyecto de Computer Vision}
    {}
    {Dec 2024 - Jan 2025}
    {
      \begin{cvitems}
         \item Desarrollo de modelos de deep learning para clasificación y segmentación de células sanguíneas (detección de leucemia linfoblástica aguda).
        \item Modelo de topología UNET logró un ROC AUC del 0.99 y F1 del 0.97 en clasificación, un ROC AUC del 0.98 y BinaryIoU del 0.94 en segmentación.
        \item Despliegue como servicio REST y aplicación web interactiva para inferencia en tiempo real.
      \end{cvitems}
      \vspace{14pt}
      \begin{itemize}[leftmargin=*]
        \small
        \item[] \textbf{Tecnologías}: TensorFlow, Keras, Streamlit, Flask, Docker, Kubernetes.
        \hspace{3em}
        \textbf{Enlaces}: 
        \href{https://github.com/aletbm/Blood_Cell_Cancer_Prediction}{Repositorio} | 
        \href{https://github.com/aletbm/Blood_Cell_Cancer_Prediction/blob/main/analysis/blood-cell-cancer-prediction.ipynb}{Análisis} | 
        \href{https://bloodcellcancerprediction.streamlit.app}{App}
      \end{itemize}
    }
    
%---------------------------------------------------------
  \cventry
    {}
    {Clasificación de tweets sobre desastres – Proyecto de Natural Language Processing}
    {}
    {Jul 2024 - Aug 2024}
    {
      \begin{cvitems}
        \item Desarrollo de un modelo de NLP para clasificar tweets sobre desastres reales a partir de un dataset de 10,000 tweets.
        \item Se alcanzó una puntuación F1 de 0.84 en la validación cruzada. El modelo híbrido de LSTM con BERT como backbone superó el rendimiento del modelo LSTM simple.
        \item Aplicación web interactiva en Streamlit para clasificación en tiempo real.
      \end{cvitems}
      \vspace{14pt}
      \begin{itemize}[leftmargin=*]
        \small
        \item[] \textbf{Tecnologías}: spaCy, TensorFlow, scikit-learn, Transformers (BERT), Streamlit.
        \hspace{0.25em}
        \textbf{Enlaces}: 
        \href{https://github.com/aletbm/NLP_with_Disaster_Tweets}{Repositorio} | 
        \href{https://github.com/aletbm/NLP_with_Disaster_Tweets/blob/main/analysis/nlp-with-disaster-tweets.ipynb}{Análisis} | 
        \href{https://disastertweetsclassification.streamlit.app}{App}
      \end{itemize}
    }

%---------------------------------------------------------
  \cventry
    {}
    {Predicción de la demanda minorista – Competencia ML Zoomcamp 2024}
    {}
    {Nov 2024 – Feb 2025}
    {
      \begin{cvitems}
        \item Desarrollo de un modelo de forecasting de demanda para múltiples series temporales de cuatro tiendas a lo largo de 25 meses.
        \item Implementación de técnicas avanzadas de feature engineering y optimización de memoria para mejorar el rendimiento del modelo.
        \item Logró la 4ᵃ posición en la competencia oficial de Kaggle del ML Zoomcamp 2024 (DataTalks.Club), alcanzando una puntuación RMSE de 8.99.
      \end{cvitems}
      \vspace{14pt}
      \begin{itemize}[leftmargin=*]
        \small
        \item[] \textbf{Tecnologías}: Python, Pandas, NumPy, scikit-learn, XGBoost, CatBoost.
        \hspace{3em}
        \textbf{Enlaces}: 
        \href{https://github.com/aletbm/Retail_Demand_Forecast_MLZoomCamp_Competition_2024}{Repositorio} | 
        \href{https://github.com/aletbm/Retail_Demand_Forecast_MLZoomCamp_Competition_2024/blob/main/top4-retail-demand-forecast-mlzc-compet-24.ipynb}{Análisis}
      \end{itemize}
    }

%---------------------------------------------------------
\end{cventries}
